\section{Background}
\label{sec:background}

\paragraph{Coverage-Guided Fuzzing.}
Coverage-guided fuzzing evaluates software by feeding mutated inputs into a target function while monitoring program execution.
Compilers such as Clang instrument basic blocks with feedback counters.
The fuzzing engine, such as libFuzzer~\cite{serebryany2016libfuzzer}, retains inputs that trigger new execution paths or edge transitions.
AddressSanitizer (ASan)~\cite{serebryany2012asan} and UndefinedBehaviorSanitizer (UBSan) detect memory safety violations, out-of-bounds array reads, and integer overflows at the exact point of failure.

\paragraph{Multi-Armed Bandit Scheduling.}
An adaptive mutator must decide which mutation operator to apply to a given input.
This decision maps directly to the classical \emph{multi-armed bandit problem}~\cite{auer2002ucb}.
Given $K$ mutation operators (arms), the scheduler balances exploitation (allocating pulls to operators that historically yielded new coverage edges) with exploration (trying less frequent operators).
The UCB1 algorithm computes an upper confidence bound score for each arm $i$:
\begin{equation}
\label{eq:ucb1}
\ArmScore{i} = \Reward{i} + C \cdot \sqrt{\frac{2 \ln \TotalPulls}{\Pulls{i}}}
\end{equation}
where $\Reward{i}$ represents the mean reward of arm $i$, $\Pulls{i}$ is the pull count of arm $i$, $\TotalPulls = \sum_{j} \Pulls{j}$ is the total number of pulls, and $C$ is an exploration parameter.
In coverage-guided fuzzing, the reward corresponds to discovering previously unseen branch edges.

\paragraph{Structured 3D Asset Formats.}
The glTF 2.0 standard~\cite{khronos_gltf2} stores 3D scenes as structured JSON metadata paired with binary buffers.
The binary GLB container consists of a 12-byte header followed by chunk entries.
Parsers such as \fastgltf decode JSON elements (nodes, meshes, accessors) and map byte offsets into binary vertex buffers.
Fuzzing such formats requires maintaining structural coherence; purely random byte flips frequently fail at early syntax verification stages.

\paragraph{Multimedia Engines.}
Game engines such as \raylib integrate graphics, audio streams, and event serialization in C.
Fuzzing these components requires simulating display contexts and audio buffers in memory without physical hardware.

\paragraph{Asynchronous Web Architectures.}
Modern C++ web servers such as \drogon handle high-concurrency HTTP traffic through asynchronous I/O.
The networking core uses thread-bound event loops that manage socket buffers and callback chains.
Testing these parsers requires driving event loops synchronously in memory to avoid operating system socket latency.
